\documentclass[11pt]{article}
\usepackage[T1]{fontenc}
\usepackage[utf8]{inputenc}
\usepackage{amsmath,amssymb,amsthm}
\usepackage[margin=1in]{geometry}
\usepackage{hyperref}
\newtheorem{theorem}{Theorem}
\title{A 34-term unit-fraction expansion with minimum denominator 18}
\author{Kyunghan Bae}
\date{11 October 2026}
\begin{document}
\maketitle
\begin{abstract}
We exhibit 34 distinct integers at least 18 whose reciprocals sum to one.
An integer certificate proves $k(18)\le34$. Together with the previously
reported computational lower bound, reproduced using the same recurrence with
additional pruning in the accompanying code, this gives $33\le k(18)\le34$.
\end{abstract}

Let $k(N)$ be the least number of distinct integers
$N\le d_1<\cdots<d_k$ such that $\sum_{i=1}^k 1/d_i=1$.

\section{The construction}
\begin{theorem}
Let
\[
\begin{split}
D=\{&18,19,20,21,22,23,24,25,26,27,28,29,30,31,\\
&32,33,34,35,36,38,39,40,43,44,45,46,\\
&48,50,52,54,63,275,14143068,10457416992\}.
\end{split}
\]
Then $|D|=34$, $\min D=18$, and $\sum_{d\in D}1/d=1$.
In particular, $k(18)\le34$.
\end{theorem}
\begin{proof}
The least common multiple of the displayed integers is
$L=6\,209\,352\,774\,424\,800$.
Exact integer addition gives
\[
\sum_{d\in D}\frac{L}{d}=6\,209\,352\,774\,424\,800=L.
\]
Divide by $L$. The cardinality, distinctness, and minimum follow directly
from the displayed set.
\end{proof}

The construction was found by replacing five terms of the published
35-term certificate in
White's report~\cite{white} with four, using
\[
\begin{split}
&\frac1{42}+\frac1{640}+\frac1{337280}
 +\frac1{1735650}+\frac1{68479214496}\\
&\hspace{12mm}=\frac1{46}+\frac1{275}
 +\frac1{14143068}+\frac1{10457416992}.
\end{split}
\]
Multiplying both sides by $M=636\,856\,694\,812\,800$ gives the same
integer $16\,160\,598\,373\,892$. The four replacement denominators
are distinct and absent from the 30 retained terms.

\section{Lower bound and the remaining case}
White's 11 August 2026 report gives $33\le k(18)\le35$, using exact
exclusions of 31 and 32 terms~\cite{white}. We reproduced the 32-term
exclusion with separately written Python exact integers and fractions,
C\# BigInteger, and C++ GMP implementations. All three completed with
814,918 visited states, 134,999 prime-condition prunes, 648,193 quadratic
terminal prunes, and 31,726 stored failed states. The programs share the
same mathematical recurrence; their agreement cross-checks implementation
and arithmetic.

For a positive reduced residual $a/b$, $t\ge2$ remaining terms, and next
denominator at least $D_0$, the first denominator lies in
\[
\max(D_0,\lfloor b/a\rfloor+1)\le x\le\lfloor tb/a\rfloor.
\]
Exact harmonic upper bounds and prime-power conditions prune impossible
branches. With two terms left, $y=bx/(ax-b)$ is checked by exact integer
division when a candidate reaches this step. No denominator cap is imposed.
In the completed 32-term run, the quadratic bound rejected every two-term
state before enumeration, so zero terminal candidates were tested. Proof details and completed
logs accompany this note as \texttt{proofs.md},
\texttt{early-n18-t32.log}, \texttt{independent-native-validation.jsonl},
and \texttt{cpp-n18-t32.log}.

Excluding 32 terms also excludes every shorter expansion. Repeatedly
replace its largest denominator $d$ using
\[
\frac1d=\frac1{d+1}+\frac1{d(d+1)}.
\]
Both new denominators exceed every retained denominator, so a shorter
expansion would extend to an impossible 32-term expansion. Hence
$k(18)\ge33$.

The existence of a 33-term expansion remains open in this work.
The asymptotic question in Erd\H{o}s problem 295,
$k(N)-(e-1)N\to+\infty$, also remains open~\cite{bloom}.

\section{A self-contained certificate check}
The following Python 3.9+ program uses only the standard library.
\begin{verbatim}
from fractions import Fraction
from math import lcm
D = [18,19,20,21,22,23,24,25,26,27,28,29,30,31,
     32,33,34,35,36,38,39,40,43,44,45,46,
     48,50,52,54,63,275,14143068,10457416992]
assert len(D) == len(set(D)) == 34 and D == sorted(D)
assert min(D) == 18
assert sum((Fraction(1,d) for d in D), Fraction()) == 1
L = lcm(*D)
assert L == 6209352774424800
assert sum(L//d for d in D) == L
old = [42,640,337280,1735650,68479214496]
new = [46,275,14143068,10457416992]
M = lcm(*old, *new)
assert M == 636856694812800
assert sum(M//d for d in old) == sum(M//d for d in new)
assert sum(M//d for d in old) == 16160598373892
print('34-term certificate verified exactly')
\end{verbatim}
The full certificate was also checked with an independently written
JavaScript BigInt calculation. Verification of the upper bound requires
only this finite identity, without running the search or installing an
optimization solver.

\section*{Computational contribution}
OpenAI Codex (GPT-6 session) agents developed the search and verification
code, found the displayed construction, performed the exact arithmetic
checks, and assisted in preparing this note.

\begin{thebibliography}{9}
\bibitem{white}
Patrick White with Claude.
\emph{Erd\H{o}s \#295: $k(17)=32$, and now $33\le k(18)\le35$
by exact exhaustion}. Erd\H{o}s Problem a Day, updated 11 August 2026.
\url{https://erdosproblemaday.com/day/295-k17-exact}.
\bibitem{bloom}
Thomas Bloom, ed. \emph{Erd\H{o}s Problem \#295}.
\url{https://www.erdosproblems.com/295}.
\end{thebibliography}
\end{document}
